AI coding agents have moved from suggesting code to acting in repositories on their own. Systems such as SWE-agent take a GitHub issue as input and try to resolve it end to end~\cite{yang2024sweagent}, and the benchmark most often used to evaluate them, SWE-bench, is itself built from real GitHub issues and the pull requests that closed them~\cite{jimenez2024swebench}. Developers adopt these tools quickly: in a controlled experiment, programmers with GitHub Copilot finished a task 55.8\% faster than those without it~\cite{peng2023copilot}, and AIDev, a recent dataset of agent activity on GitHub, contains hundreds of thousands of pull requests opened by agents~\cite{li2026aidev}. The issue tracker sits at the start of this workflow: it is where work is requested, discussed and planned, and its use in open-source projects has been studied at scale~\cite{bissyande2013issues}.

Agents now also appear \emph{inside} issue trackers, and they do so in different roles. A review bot such as CodeRabbit opens follow-up issues for problems it found in a pull request. Sweep is summoned by a human who adds a label or a title prefix. Other agents show up only because someone mentions or discusses them. These roles matter to maintainers. When projects adopted earlier code-review bots, their pull-request activity changed: more pull requests were merged and developers communicated less~\cite{wessel2020effects}, and developers describe unwanted bot activity as noise that distracts them from their work~\cite{wessel2021disturb}. For researchers, the roles create a measurement problem: datasets attribute issues to agents with heuristic signals, and the most recent of these datasets warns that its agent label is a candidate attribution, not a verified one~\cite{ghaleb2026codage}. If ``issues attributed to agent X'' mixes issues the agent wrote with issues that merely mention it, comparisons between agents may say more about the attribution than about the agents.

This paper compares the GitHub issues attributed to \NAgents{} AI coding agents. Following the Goal--Question--Metric approach~\cite{basili1994gqm}, it asks three research questions:
\begin{itemize}
\item \RQ{1}: \emph{How do issues attributed to different agents differ in volume, state and labelling?} Before comparing speed or content, maintainers need a baseline of what each agent's issues look like and how often they are closed.
\item \RQ{2}: \emph{Does the time-to-close of these issues differ across agents, and by how much?} Time-to-close is a common proxy for how quickly maintainers act on an issue~\cite{kikas2016lifetime}, but a third of the issues in the data are still open, so the answer depends on how open issues are treated.
\item \RQ{3}: \emph{Which topics do the issues raise, and does that differ by agent?} Topics show what kind of work an agent's issues represent, which helps explain the differences found in RQ1 and RQ2.
\end{itemize}

Most empirical work on coding agents studies their pull requests: how often they are accepted~\cite{li2025aiteammates,pinna2026comparing}, what tasks developers hand to them~\cite{watanabe2025agentic}, and how their changes and descriptions differ from those of humans~\cite{pham2026alignment}. Closer to the setting of this study, Ghaleb showed that agents leave behavioural fingerprints in their pull requests~\cite{ghaleb2026fingerprinting}, and Selvanayagam and Ghaleb studied agents that review other agents' code~\cite{selvanayagam2026aitoai}. Sayagh asked which issues make good tasks for an agent~\cite{sayagh2025copilot}, but for a single agent. To the best of the researcher's knowledge, no study has compared the issues attributed to many agents on their characteristics, their resolution time and their topics together.

The study covers \NIssues{} issues from \NRepos{} repositories, attributed to \NAgents{} agents. The data turned out to be less uniform than the label suggests. For CodeRabbit and Devin, most issues were written by the agent's own bot account (\OwnBotShareCoderabbit\% and \OwnBotShareDevin\%), while for Claude Code, Sweep and Jules almost none were. This role is associated with the outcome: within the same agent, bot-written issues are closed less often than the others, by 18 to 24 percentage points. How open issues are handled also matters. Ranking the agents by the time-to-close of their closed issues gives a very different order from a survival analysis that counts the issues still open (Spearman \(\rho = \RankRho\)). The topics confirm the picture from a different angle. CodeRabbit's own issues are follow-ups from pull-request reviews, with code-level fixes as their largest topic, and this result holds under two different topic models, even though the two models divide the corpus differently.

The paper contributes the following:
\begin{enumerate}
\item a comparison of issues attributed to \NAgents{} AI coding agents on volume, state, labelling, time-to-close and topics;
\item evidence that the role an agent plays in an issue, written by its bot, triggered by a human or merely mentioned, differs sharply between agents and is associated with whether the issue is closed;
\item a demonstration that comparing agents by the closed issues alone can reverse their ranking, and that survival analysis avoids this;
\item a manual, blind validation of the attribution, the topic labels and the language flag on random samples, and a replication package with all scripts, independent verification checks and the validation answers.
\end{enumerate}

\textbf{Paper organisation.} Section~\ref{sec:related} reviews related work. Section~\ref{sec:method} describes the data and the methodology. Sections~\ref{sec:rq1} to~\ref{sec:rq3} answer the three research questions. Section~\ref{sec:discussion} discusses implications, Section~\ref{sec:threats} the threats to validity, and Section~\ref{sec:conclusion} concludes.
